\documentclass[10pt,a4paper]{amsart}
\usepackage[margin=19mm]{geometry}
\usepackage{amsmath,amssymb,mathtools}
\usepackage[scr=rsfs,cal=euler]{mathalfa}
\usepackage{xfrac}
\usepackage[unicode,hidelinks,bookmarks=false]{hyperref}
\newcommand{\ie}{i.e.\ }
\newcommand{\cf}{cf.\ }
\newcommand{\resp}{respectively\ }
\newcommand{\Subsection}{\subsection}
\providecommand{\nobreakdash}{\nobreak\hbox{-}}
\setlength{\emergencystretch}{2em}
\multlinegap0pt
\usepackage{bm}
\usepackage{enumerate}
\usepackage{amssymb}
\usepackage{xcolor}
\newtheorem{theorem}{Theorem}
\newtheorem{cor}{Corollary}
\newtheorem{lemma}{Lemma}
\newcommand{\BH}{\mathcal{B(H)}}
\newcommand{\Conj}{\operatorname{Conj}(\cH)}
\newcommand{\SC}{\mathcal{S}_C}
\newcommand{\Span}{\operatorname{span}}
\newcommand{\cH}{{\mathcal{H}}}
\title{Linear maps preserving $C$-symmetry operators}
\author{Mayu Anzai, Shiho Oi}
\date{Source: arXiv:2608.24583v2 (CC BY 4.0)}
\begin{document}
\maketitle

\noindent\emph{Source and licence.} Linear maps preserving $C$-symmetry operators. Version arXiv:2608.24583v2 (CC BY 4.0), distributed by arXiv under the Creative Commons Attribution 4.0 International licence, http://creativecommons.org/licenses/by/4.0/, as recorded in the arXiv licence metadata for this exact version and shown on the abstract page https://arxiv.org/abs/2608.24583v2. Full licence notice: \texttt{licenses/arxiv-2608.24583-CC-BY-4.0.txt}.\par\smallskip

\noindent\textit{Editorial scope.} Counted unit: the article's theorem main0 (source0.tex lines 190--209). The article's complete original proof of it is delivered with it (source0.tex lines 599--664, one \texttt{proof} environment of 66 lines, which the article places away from the statement). No mathematics is written, repaired or re-derived: the statement and the proof are the article's own lines, in the article's own notation, and no retained line is substituted or re-flowed.  Retained uncounted as the source-local context the proof needs: lines 432--437; lines 439--443; lines 491--493; lines 513--516; lines 566--571.  Nothing was omitted from any retained span, so the package carries no omission record.  No retained line contains a citation, so the wrapper carries no bibliography and no bibliography entry.  No retained line contains a figure environment, an image-inclusion command or drawing code, so no image is delivered and the package declares no asset dependency.  Editorial formatting only, in addition: an AMS \texttt{amsart} preamble that restates the article's own declarations and macros used by the retained lines, namely the environments \texttt{theorem}, \texttt{cor}, \texttt{lemma} on separate counters, together with the standard packages those lines need.  The retained environments print in this document as Theorem 1, Corollary 1, Lemma 1 in order of appearance, whereas the article prints the same items with the numbering of its own full document.  The complete native source of the article as distributed in the arXiv e-print for this exact version is bundled unchanged as \texttt{sources/208491/source0.tex}.
\begin{cor}\label{dim2o}
Let $x \in \cH$ with $\|x\|=1$. Then 
\[
\bigcap_{C \in \Lambda_x} \SC=\Span\{x\otimes x, I\}.
\]
\end{cor}

\begin{lemma}\label{dim2t}
Let $\Lambda \subset \Conj$. Then  $\dim (\bigcap_{C \in \Lambda} \SC)=2$  if and only if 
there is a projection $P \in \BH$ with $P \neq I, 0$ such that %$C(P(\cH)) =P(\cH)$ for any $C \in \Lambda$ and 
$\bigcap_{C \in \Lambda} \SC=\Span \{P, I\}$.
\end{lemma}

\begin{lemma}\label{I3}
Let $\cH$ be a separable Hilbert space. Let $T: \BH \to \BH$ be a continuous bijective linear map  and  $\varphi: \mathcal{E}(\cH) \to \mathcal{E}(\cH)$ be a bijection satisfying $T(\mathcal{D}_{\{e_n\}}) = \mathcal{D}_{\varphi(\{e_n\})}$. There is $c \in \mathbb{C} \setminus \{0\} $ such that $T(I)=cI$.
\end{lemma}

\begin{lemma}\label{lift}
Suppose that $\mathcal{F}:\mathcal{E}(\cH)/ \sim \to \mathcal{E}(\cH)/ \sim$
is a bijection. Then there is a bijection $\varphi:\mathcal{E}(\cH) \to \mathcal{E}(\cH)$ such that $q\circ \varphi=\mathcal{F} \circ q$.
\end{lemma}

\begin{lemma}\label{plane}
We have
\begin{multline*}
    \{\SC \mid C \in \Conj\}\\=\{\Span\{I, P\cdot\sigma\} \mid P \text{ is  a two-dimensional subspace of $\mathbb{R}^3$}\}
\end{multline*}
\end{lemma}

\begin{theorem}\label{main0}
Let $\cH$ be a complex Hilbert space with $\dim \cH=2$.  Let $T: M_2(\mathbb{C}) \to M_2(\mathbb{C})$ be a bijective linear map. Then the following are equivalent.
\begin{itemize}
\item[(1)] there is a bijection $\psi: \Conj \to \Conj$ which satisfies  
 \[
T(\SC) = {\mathcal{S}_{\psi(C)}},
\]
\item[(2)] there is a bijection $\varphi: \mathcal{E}(\cH) \to \mathcal{E}(\cH)$,  which satisfies  
\[
T(\mathcal{D}_{\{e_n\}}) = \mathcal{D}_{\varphi(\{e_n\})},
\] 
\item[(3)] there exists $c, \lambda \in \mathbb{C} \setminus \{0\}$, a linear functional $\ell$ on $\mathbb{C}^3$, and $B \in GL(3,\mathbb{R})$ such that 

\[
T(aI+x\cdot \sigma)=c aI+\ell(x)I+\lambda Bx
\] 
for  $a \in \mathbb{C}$ and $x \in \mathbb{C}^3$.
\end{itemize}

\end{theorem}

\begin{proof}[{\textbf {Proof of Theorem \ref{main0}}}] 
$\bm{(1) \Rightarrow (2)}.$ Let $\{e_1,e_2\} \in \mathcal{E}(\cH)$ and  $C \in \Lambda_{e_1}$. Corollary \ref{dim2o} shows that 
\[
\bigcap_{C \in \Lambda_{e_1}} \SC=\Span\{e_1 \otimes e_1, I \}=D_{\{e_1,e_2\}}. 
\]
Since $T(\bigcap_{C \in \Lambda_{e_1}} \SC)=\bigcap_{C \in \Lambda_{e_1}} \mathcal{S}_{\psi(C)}$ and $\dim \bigcap_{C \in \Lambda_{e_1}} \SC=2$, we have $\dim \bigcap_{C \in \Lambda_{e_1}} \mathcal{S}_{\psi(C)}=2$. By Lemma \ref{dim2t},  there is $e'_1 \in \cH$ with $\|e'_1\|=1$ such that 
\[
\bigcap_{C \in \Lambda_{e_1}} \mathcal{S}_{\psi(C)}=\Span\{e'_1 \otimes e'_1, I \}=D_{\{e'_1,e'_2\}},
\]
where $\{e'_1, e'_2\} \in \mathcal{E}(\cH)$. 
Hence, we have 
\begin{equation}\label{diago}
T(D_{\{e_1,e_2\}})=D_{\{e'_1,e'_2\}}, 
\end{equation}
for any $\{e_1,e_2\}  \in \mathcal{E}(\cH)$. Since a similar argument holds for $T^{-1}$, there is a bijective map $\mathcal{F}: \mathcal{E}(\cH)/ \sim \to \mathcal{E}(\cH)/ \sim$ such that 
\[
T(D_{[e_n]})=D_{\mathcal{F}([e_n])}.
\]
Applying Lemma \ref{lift}, there is a bijection $\varphi$ such that $q\circ \varphi=\mathcal{F} \circ q$. Therefore we have
\[
T(D_{\{e_n\}})=T(D_{[e_n]})=D_{\mathcal{F}([e_n])}=D_{[\varphi(\{e_n\})]}=D_{\varphi(\{e_n\})}.
\] 
This completes the proof of  $(1)\Rightarrow(2)$. \\
%As $x$ was chosen arbitrarily, \eqref{diago} 
%%%%%%%%%%%%%%%%%%%%%%%%%%%%%%%%%%%%%%
$\bm{(2) \Rightarrow (3)}.$ 
By Lemma \ref{I3}, there is $c \in \mathbb{C} \setminus \{0\}$ such that $T(I)=cI$. Since $T$ is a complex linear map, there are linear maps $\ell: \mathbb{C}^3 \to \mathbb{C}$ and $L: \mathbb{C}^3 \to \mathbb{C}^3$ such that 
\[
T(x \cdot \sigma)=\ell(x)I+L(x)\cdot \sigma.
\]
Since $T$ is a bijection, $L$ is an invertible matrix. For any $y=(y_1, y_2, y_3) \in \mathbb{R}^3 \setminus \{0\}$, there is $\{e_1, e_2\} \in \mathcal{E}(\cH)$ such that $\frac{1}{2}(I+\frac{y}{\|y\|}\cdot \sigma)=e_1\otimes e_1$. Then %For any rank one projection $e\otimes e$, there is a rank one projection $f \otimes f$ such that 
\[
T(\Span\{e_1\otimes e_1, I\})=\Span\{f_1 \otimes f_1, I\},
\]
where $\varphi(\{e_1, e_2\})=\{f_1, f_2\}$. Thus there is $z=(z_1, z_2, z_3) \in \mathbb{R}^3$ with $z_1^2+z_2^2+z_3^2=1$ such that $f_1 \otimes f_1=\frac{1}{2}(I+z \cdot \sigma)$. 
It implies 
\begin{multline*}
\Span\{I, L(y)\cdot \sigma\}=T(\Span\{I, y/\|y\| \cdot \sigma\})\\=T(\Span\{e_1\otimes e_1, I\})=\Span\{f_1 \otimes f_1, I\}=\Span\{I, z \cdot \sigma\}.
\end{multline*}
As $\{I, \sigma_1, \sigma_2, \sigma_3\}$ are linearly independent,  we have $L(y)=\lambda z$ for some $\lambda \in \mathbb{C}$. Let $e_1=(1,0,0)$ , $e_2=(0,1,0)$ and $e_3=(0,0,1)$ be the standard basis on $\mathbb{R}^3$. Then there are $b_i \in \mathbb{R}^3$ with $\|b_i\|=1$ such that $L(e_i)=\lambda_i b_i$ for some $\lambda_i \in \mathbb{C}$ for $i=1,2,3$. 
In addition as $e_1+e_2+e_3 \in \mathbb{R}^3$, there is $b \in \mathbb{R}^3$ such that $L(e_1+e_2+e_3)=\lambda b$ for some $\lambda \in \mathbb{C}\setminus\{0\}$. Since $L$  is an invertible matrix and $\{e_1, e_2, e_3\}$  is an orthonormal basis, there is $\{b_1, b_2, b_3\}$ is linearly independent. As $b_i \in \mathbb{R}^3$ and $b \in \mathbb{R}^3$, there are $r_i \in \mathbb{R}$ such that $b=r_1b_1+r_2b_2+r_3b_3$. It implies that
\[
\lambda_1 b_1+\lambda_2 b_2+\lambda_3 b_3=L(e_1+e_2+e_3)=\lambda b= \lambda (r_1b_1+r_2b_2+r_3b_3).
\]
We get $\lambda_i=\lambda r_i$ for any $i=1,2,3$. We define $B \in GL(3,\mathbb{R})$ by $Be_i=r_ib_i$ for any $i=1,2,3$. Then we have $L=\lambda B$. Therefore we have
\[
T(aI+x\cdot \sigma)=c aI+\ell(x)I+\lambda Bx.
\] 

$\bm{(3) \Rightarrow (1)}.$
Applying Lemma \ref{plane}, for any $C \in \Conj$, we have 
\[
\SC=\Span\{I, P\cdot\sigma\},  
\]
where $P$ is a  two-dimensional subspace of $\mathbb{R}^3$. For any $x \in P$, we have
\[
T(aI+x \cdot \sigma)=c aI+\ell(x)I+\lambda Bx \cdot \sigma \in \Span\{ I, BP \cdot \sigma\}.
\]
Since $B \in GL(3, \mathbb{R})$, we get $BP$  is a  two-dimensional subspace of $\mathbb{R}^3$. Hence there is $D \in \Conj$ such that $\Span\{ I, BP \cdot \sigma\}=\mathcal{S}_D$. As $T$ is a bijective linear map, we obtain $T(\SC)=\mathcal{S}_D$. The same argument applies to $T^{-1}$. Thus, by choosing a section of the quotient space of conjugations modulo the relation $C\sim e^{i\theta}C$, we can define, as above, a bijection $\psi:\Conj \to \Conj$ satisfying
\[
T(S_C)=S_{\psi(C)}.
\]

%%%%%%%%%%%%%%%%%%%%%%%%%%%%%%%%%%%%%%%%%%%

\end{proof}

\end{document}
